% \jz{add the Big 5 Personalities paper, LLM economicus paper, and the Game-theoretic LLMs paper?}
% \lh{If I was thinking about what to add here, I'd try to think about key themes and ideas rather than a random assortment of papers that seem somewhat related. Probably the key themes in this regard are cooperative capabilities and dispositions? Maybe also general stuff on strategic reasoning in LLMs.

% Some ideas of relevant papers (not exhaustive):
% \begin{itemize}
% \item Cooperation and Control in Delegation Games
% \item Quantifying Misalignment Between Agents
% \item Quantifying Differences in Reward Functions
% \item On Measuring Social Intelligence: Experiments on Competition and Cooperation
% \item Beyond the high score: Prosocial ability profiles of multi-agent populations
% \item  Rating the skill of synthetic agents
% in competitive multi-agent environment
% \item Towards the Scalable Evaluation of Cooperativeness in Language Models
% \item  Strategic Reasoning with Language Models
% \item  Decoupling Strategy and Generation
% in Negotiation Dialogues
% \end{itemize}
% I didn't add many papers on the dispositions of LLMs here, but would recommend citing a few on that. Maybe some of Dan Hendrycks' work (such as the ETHICS or MACHIAVELLI datasets)?

% Probably also good to briefly mention earlier work on automated negotiation in AI (i.e., not necessarily using LLMs). That keyword (`automated negotiation') should pull up a bunch of relevant papers, and there is also a bunch of relevant MARL work in that domain.
% }

LLM negotiation has quickly become a useful testbed for studying strategic reasoning in natural language.
Early work introduced benchmark frameworks and capability decompositions for negotiation \citep{bianchi_how_2024,kwon-etal-2024-llms}, and subsequent papers have explored bargaining in a range of simulated settings, including buyer--seller exchange, resource trade, and multi-issue stakeholder negotiation \citep{abdelnabi_cooperation_2024,davidson_evaluating_2024,chatterjee_agreemate_2024,noh_llms_2024}.
However, more capable models do not necessarily become rational or unbiased negotiators.
A recent frontier-model re-evaluation of NegotiationArena finds that stronger models still diverge into model-specific strategic signatures rather than converging to a common bargaining policy, while also exhibiting persistent numerical and semantic anchoring effects and clear dominance patterns in which some models systematically extract higher payoffs than others \citep{rios_illusion_2025}.
These findings are closely related to our own concern that general-purpose capability does not map monotonically onto strategic performance.

A complementary line of work studies exploitability under adaptive strategic pressure rather than average performance against a fixed pool of opponents.
\citet{wang_profit_2026} introduce \emph{profit-driven red teaming}, in which an opponent is optimized directly for profit using only scalar outcome feedback in four auditable economic games: ultimatum bargaining, first-price auctions, bilateral trade, and a provision-point public goods game.
They show that agents which appear strong against static baselines can become reliably exploitable under adaptive pressure, with the learned opponent discovering tactics such as probing, anchoring, and deceptive commitments.
This is complementary to our setting: instead of red-teaming a fixed target, we study how bargaining outcomes vary across model capabilities and across systematically controlled cooperation--competition regimes.

Our work also connects to research on scaling and adaptation in strategic interaction.
Prior scaling-law work on strategic behavior has mostly focused on \emph{multi-agent reinforcement learning} (MARL) settings such as Hex, Connect Four, and Pentago, rather than natural-language bargaining between LLM agents \citep{jones_scaling_2021,neumann_scaling_2022}.
More recently, \citet{liu_scaling_2026} show that repeated negotiation performance can be improved at inference time by allocating additional computation to opponent modeling and simulated best-response search, without updating model parameters.
Their results suggest that strategic performance depends not only on base model capability, but also on how test-time computation is spent.
More broadly, scaling laws have been studied for single-agent RL \citep{hilton_scaling_2023}, vision \citep{Zhai2022}, LLMs \citep{hoffmann_training_2022,kaplan_scaling_2020}, and scalable oversight \citep{engels_scaling_2025}; our paper brings this perspective to negotiation environments with tunable strategic structure.

Finally, recent work has begun to push beyond bilateral outcome bargaining toward richer multi-party negotiation settings.
\citet{benac_benchmark_2026} introduce a benchmark for multi-party, sequential action-level negotiation with binding commitments and terminal-only rewards, together with a configurable generator and document-grounded climate negotiation instances derived from real negotiation artifacts.
Their results emphasize that performance is strongly regime-dependent and that simple zero-shot LLM negotiation remains uncompetitive in these large, path-dependent settings.
This is especially relevant to our planned extensions beyond two-agent bargaining.
Whereas their emphasis is on benchmark design and value-function approximations for long-horizon commitment games, our present paper focuses on autonomous LLM negotiation in controlled environments with tunable cooperation--competition structure and broad capability sweeps.

% \lh{As a heuristic, the related work section for a paper this long should be at least half a page.}
